\subsection{分次理想的维数与链}

在本节中，记 dimgr(H) 为 H 的分次素理想链长度的上确界；同样，若 p 是 H 的分次素理想，记 htgr(p) 为 H 的、以 p 为最大元的分次素理想链长度的上确界。若 p 是分次素理想，则有 \(p \cap S = \varnothing\)；否则，p 将包含一个次数 0 分量等于 1 的元素，由于 p 是分次理想，遂包含 1。因而，从 H 的分次素理想集合到 \(p \mapsto S^{-1}p\) 的素理想集合的映射 \(S^{-1}H\) 是单射且递增的（II, § 2, no. 5, 命题 11）；因此，结合 § 1, no. 3 的命题 6 及命题 7 的推论，有：

命题 2. --- a) 有 \(\dim \mathrm{gr}(\mathrm{H}) \leqslant \dim (\mathrm{S}^{-1}\mathrm{H}) \leqslant \dim (\mathrm{H})\)。

\begin{enumerate}
\def\labelenumi{\alph{enumi})}
\setcounter{enumi}{1}
\tightlist
\item
  对于 H 的每个分次素理想 \(\mathfrak{p}\)，有 \(\operatorname{htgr}(\mathfrak{p}) \leqslant \operatorname{ht}(\mathrm{S}^{-1}\mathfrak{p}) = \operatorname{ht}(\mathfrak{p})\)。
\end{enumerate}

对于 H 的每个理想 a，记 \(a^{gr}\) 为包含于 a 的最大分次理想；有 \(a^{gr} = \sum_{n} (a \cap H_n)\)。

引理 1. --- a) 若 p 是 H 的素理想，则 \(\mathfrak{p}^{\mathrm{gr}}\) 是素理想。

\begin{enumerate}
\def\labelenumi{\alph{enumi})}
\setcounter{enumi}{1}
\item
  H 的分次理想集合中每个不同于 H 的极大元，都是包含 \(H_{\geq1}\) 的 H 的极大理想。
\item
  H 的每个极小素理想都是分次的。
\item
  这由 III, § 1, no. 4 的命题 4 得出。
\item
  设 m 是 H 的不同于 H 的分次理想。则有
\end{enumerate}

\[
\mathfrak {m} \subset (\mathfrak {m} \cap \mathrm{H} _ {0}) + \mathrm{H} _ {\geqslant 1} \neq \mathrm{H}.
\]

如果 m 是极大理想，则有 \(m = m_{0} + H_{\geq 1}\)，其中 \(m_{0}\) 是 \(H_{0}\) 的极大理想，从而得到 b)。

\begin{enumerate}
\def\labelenumi{\alph{enumi})}
\setcounter{enumi}{2}
\tightlist
\item
  设 p 是 H 的极小素理想。由于 \(p^{gr}\) 由 a) 知是素理想且包含于 p，故有 \(p = p^{gr}\)，从而得到 c)。
\end{enumerate}

引理 2.--- 设 p、q 是 H 的素理想，满足 q ⊂ p 且 q ≠ p。若 q \(^{gr}\) = p \(^{gr}\)，则 q 是分次的，p 不是分次的，且 ht(p/q) = 1。

注 1.--- 沿用第 1 号例 1 的记号。引理 2 表明：若 \(\mathbf{C}^{n+1}\) 中两个不可约子簇 Y、Z 具有相同的投影锥，且 \(Z \subset Y\)、\(Z \neq Y\)，则 Y 是 Z 的投影锥，而 Z 在 Y 中的余维数为 1。*

以 \(\mathrm{H} / \mathfrak{q}^{\mathrm{gr}}\) 替换 H，将问题化为 \(\mathfrak{q}^{\mathrm{gr}} = \{0\}\) 的情形。此时 H 是整环（引理 1, a)），\(\mathfrak{p}^{\mathrm{gr}} = 0\)，而要证明的是 \(\operatorname{ht}(\mathfrak{p}) \leqslant 1: this will indeed imply that  \operatorname{ht}(\mathfrak{q}) = 0\)，从而有 \(\mathfrak{q} = \{0\}\)。由于 \(\mathfrak{p}^{\mathrm{gr}} = \{0\}\)，对于每个 \(\mathfrak{p} \cap \mathrm{H}_n = \{0\}\)，有 \(n\)，且 \(\mathfrak{p}\) 与乘法子集 \(\mathrm{T} = \bigcup_{n} (\mathrm{H}_n - \{0\})\) 不相交。环 \(\mathrm{H}_{\mathfrak{p}}\) 因而同构于 \(\mathrm{T}^{-1}\mathrm{H}\) 的一个分式环，从而有

\[
\mathrm{ht} (\mathfrak {p}) = \dim (\mathrm{H} _ {\mathfrak {p}}) \leqslant \dim (\mathrm{T} ^ {- 1} \mathrm{H})
\]

（§ 1, no. 3, 命题 6 和命题 7）。但是，根据 V, § 1, no. 8 的引理 4，T \(^{-1}\) H 是域，或同构于环 K{[}X, X \(^{-1}\) {]}，其中 K 是域；因而有 dim(T \(^{-1}\) H) ≤ 1 及 ht(p) ≤ 1，这正是要证明的结论。

命题 3. --- 设 p 是 H 的素理想。若 p ≠ p \(^{gr}\)，则有 ht(p \(^{gr}\)) = ht(p) - 1。根据引理 1 a)，理想 p \(^{gr}\) 是素理想且包含于 p，故 ht(p \(^{gr}\)) ≤ ht(p) - 1。当 ht(p \(^{gr}\)) = +∞ 时命题显然，因此可设 ht(p \(^{gr}\)) \textless{} +∞。我们按 ht(p \(^{gr}\)) 归纳证明不等式 ht(p) ≤ ht(p \(^{gr}\)) + 1。只需证明：对于包含于 p 且不同于 p 的每个素理想 q，有 ht(q) ≤ ht(p \(^{gr}\))。按 q \(^{gr}\) ≠ p \(^{gr}\) 或 q \(^{gr}\) = p \(^{gr}\) 分两种情形。若 q \(^{gr}\) ≠ p \(^{gr}\)，则 ht(q \(^{gr}\)) \textless{} ht(p \(^{gr}\))；我们有

\[
\mathrm{ht} (\mathfrak {q}) \leqslant \mathrm{ht} (\mathfrak {q} ^ {\mathrm{gr}}) + 1,
\]

根据归纳假设（当 q ≠ q \(^{gr}\) 时）以及显然性（当 q = q \(^{gr}\) 时），我们得到 ht(q) ≤ ht(q \(^{gr}\)) + 1 ≤ ht(p \(^{gr}\))，这正是要证明的结论。若 q \(^{gr}\) = p \(^{gr}\)，则根据引理 2 有 q = q \(^{gr}\)，故 ht(q) = ht(q \(^{gr}\)) ≤ ht(p \(^{gr}\))，再次得到所需结论。

定理 1. --- 假设 H 是诺特环。

\begin{enumerate}
\def\labelenumi{\alph{enumi})}
\item
  H 的每条分次素理想链，若作为分次素理想链是饱和的，则作为素理想链也是饱和的。
\item
  对于 H 的每个分次素理想 p，有 htgr(p) = ht(S \(^{-1}\) p) = ht(p)。
\item
  我们有 \(\dim \mathrm{gr}(\mathbf{H}) = \dim (\mathbf{S}^{-1}\mathbf{H}) = \dim (\mathbf{H}).\)
\end{enumerate}

为证明 \(a\)，只需证明：若 \(\mathfrak{p}\) 和 \(\mathfrak{q}\) 是 H 的两个不同的分次素理想，满足 \(\mathfrak{q} \subset \mathfrak{p}\)，且 \(\mathfrak{q}\) 与 \(\mathfrak{p}\) 之间的每个分次素理想都等于 \(\mathfrak{p}\) 或 \(\mathfrak{q}\)，则 \(\mathrm{ht}(\mathfrak{p}/\mathfrak{q}) = 1\)。将 \(\mathfrak{q}\) 除去后，问题化为 \(\mathfrak{q} = \{0\}\) 的情形。因此还需证明：若 H 是整环，且 \(\mathfrak{p}\) 是 H 的一个理想，在不同于 \(\neq \{0\}\) 的分次素理想中为极小元，则有 \(\mathrm{ht}(\mathfrak{p}) = 1\)。现在设 \(a\) 是 \(\mathfrak{p}\) 中的非零齐次元素，令 r 为 H 的一个素理想，使得 \(a \in \mathfrak{r} \subset \mathfrak{p}\)，并且对于这些性质是极小的（II, § 2, no. 6, 引理 2）。由于 \(\mathfrak{r}^{\mathrm{gr}}\) 是素理想（引理 1，\(a\)），且非零，故有 \(\mathfrak{r}^{\mathrm{gr}} = \mathfrak{p}\)，从而 \(\mathfrak{p} = \mathfrak{r}\)。由于 H 是整环且为诺特环，\(\mathfrak{p}\) 的高度为 1（§ 3, no. 1, 命题 1），从而得到 \(a\)。

证明 b)。设 p 是 H 的一个分次素理想。有

\[
\operatorname{htgr} (\mathfrak {p}) \leqslant \operatorname{ht} \left(\mathrm{S} ^ {- 1} \mathfrak {p}\right) \leqslant \operatorname{ht} (\mathfrak {p})
\]

（命题 2）；我们按 ht(p) ≤ htgr(p) 归纳证明不等式。若 htgr(p) = 0，则 p 在分次素理想中是极小的，因而是极小的（引理 1，c)），且 ht(p) = 0。假设 htgr(p) \textgreater{} 0，并证明：对于每个包含于 p 且不同于 p 的素理想 q，有 ht(q) ≤ htgr(p) - 1。分两种情形。若 q 是分次的，则由归纳假设得出结论。若 q 不是分次的，则有 q \(^{gr}\) ≠ p，故根据归纳假设 ht(q \(^{gr}\) ) ≤ htgr(q \(^{gr}\) )，从而根据命题 3，ht(q) ≤ htgr(q \(^{gr}\) ) + 1；还需证明 htgr(q \(^{gr}\) ) ≤ htgr(p) - 2；但若 htgr(q \(^{gr}\) ) = htgr(p) - 1，则链 q \(^{gr}\) ⊂ p 根据 a) 将是饱和的，而这不可能，因为 q \(^{gr}\) ≠ q ≠ p。

最后证明 c)。有 dimgr(H) ≤ dim(S⁻¹H) ≤ dim(H)（命题 2），还需证明 dim(H) ≤ dimgr(H)，或等价地，对 H 的每个素理想 p，有 ht(p) ≤ dimgr(H)。于是设 p 是 H 的素理想。若 p 分次，则 ht(p) = htgr(p) ≤ dimgr(H)。若 p 非分次，则根据命题 3，有 ht(p) = htgr(p\(^{\mathrm{gr}}\)) + 1；令 m 为 H 中包含 p\(^{\mathrm{gr}}\) 的极大分次理想；根据引理 1 b)，m 是极大的，因而不同于 p\(^{\mathrm{gr}}\)，且有 htgr(p\(^{\mathrm{gr}}\)) + 1 ≤ htgr(m) ≤ dimgr(H)，从而再次有 ht(p) ≤ dimgr(H)。证明完毕。

注记 2. --- 存在非诺特的分次环 H，使 dimgr(H) \textless{} dim(H)（p.~99，练习 1）。

